\documentclass[a4paper,11pt]{book}
\usepackage{listings}
\usepackage[utf8]{inputenc}
\usepackage{titlesec}
\usepackage{fancyhdr}
\usepackage[spanish,es-tabla]{babel}
\usepackage[hidelinks]{hyperref}
\usepackage{xcolor}
\usepackage{pdfpages}
\usepackage{url}
\usepackage{booktabs}
\usepackage[export]{adjustbox}
\usepackage{fancybox}

\usepackage{textcomp}

\usepackage{wrapfig}


\usepackage{float}

\usepackage{booktabs}

\usepackage{rotating}

% Información reutilizable
\newcommand{\asunto}{Cloud Computing}
\newcommand{\titulo}{Servicio de	almacenamiento	en	
cloud	de	alta	disponibilidad}
\newcommand{\grado}{Máster Profesional en Ingeniería Informática}
\newcommand{\autor}{Ernesto Serrano Collado}
\newcommand{\email}{info@ernesto.es}
\newcommand{\profesor}{Manuel Parra}
\newcommand{\escuela}{Escuela Técnica Superior de Ingenierías Informática y de Telecomunicación}
\newcommand{\universidad}{Universidad de Granada}
\newcommand{\ciudad}{Granada}
\newcommand{\vers}{Versión 0.1}
\providecommand{\keywords}{cloud computing, openstack, docker, owncloud, nginx, postgresql, ldap}

% Información archivo
\hypersetup{
	pdfauthor = {\autor\ (\email)},
	pdftitle = {\titulo},
	pdfsubject = {\asunto},
	pdfkeywords = {\keywords},
	pdfcreator = {MacTeX con el paquete TeX Live},
	pdfproducer = {pdflatex}
}

% Estilo de cabeceras
\pagestyle{fancy}
\fancyhf{}
\fancyhead[LO]{\leftmark}
\fancyhead[RE]{\rightmark}
\fancyhead[RO,LE]{\textbf{\thepage}}
\setlength{\headheight}{1.5\headheight}

% Redefinición de comandos
\renewcommand{\lstlistingname}{Fragmento de código}
\renewcommand{\lstlistlistingname}{Índice de fragmentos de código}
\renewcommand{\chaptermark}[1]{\markboth{\textbf{#1}}{}}
\renewcommand{\sectionmark}[1]{\markright{\textbf{\thesection. #1}}}

% Definición de colores
\definecolor{gray97}{gray}{.97}
\definecolor{gray75}{gray}{.75}
\definecolor{gray45}{gray}{.45}
\definecolor{gray30}{gray}{.94}
\definecolor{lightgray}{rgb}{.9,.9,.9}
\definecolor{darkgray}{rgb}{.4,.4,.4}
\definecolor{purple}{rgb}{0.65, 0.12, 0.82}
\definecolor{background}{HTML}{EEEEEE}
\definecolor{delim}{RGB}{20,105,176}
\colorlet{punct}{red!60!black}
\colorlet{numb}{magenta!60!black}

	\definecolor{dkgreen}{rgb}{0,0.6,0}
	\definecolor{gray}{rgb}{0.5,0.5,0.5}
	\definecolor{mauve}{rgb}{0.58,0,0.82}

% Listados
\lstset{
	aboveskip=0.5cm,
	backgroundcolor=\color{gray97},
	basicstyle=\scriptsize\ttfamily,
	breaklines=true,
	%commentstyle=\color{gray45},
	frame=Ltb,
	framerule=0.5pt,
	framesep=0pt,
	framexbottommargin=3pt,
	framexleftmargin=0.1cm,
	framextopmargin=3pt,
	%keywordstyle=\bfseries,
	numberfirstline = false,
	numbers=left,
	numbersep=6pt,
	%numberstyle=\tiny,
	rulesep=.4pt,
	rulesepcolor=\color{black},
	showstringspaces = false,
	%stringstyle=\ttfamily,
	  numberstyle=\tiny\color{gray},
	  keywordstyle=\color{blue},
	  commentstyle=\color{dkgreen},
	  stringstyle=\color{mauve},
	literate={á}{{\'a}}1
	         {é}{{\'e}}1
	         {í}{{\'i}}1
	         {ó}{{\'o}}1
	         {ú}{{\'u}}1
	         {ñ}{{\~n}}1
}


% Minimizar fragmentado de listados
\lstnewenvironment{listing}[1][]
	{\lstset{#1}\pagebreak[0]}{\pagebreak[0]}

% Listado definido para JavaScript
% http://tex.stackexchange.com/questions/89574/language-option-supported-in-listings/89576#89576
\lstdefinelanguage{javascript}{
	backgroundcolor=\color{background},
	basicstyle=\footnotesize,
	breaklines=true,
	captionpos=b,
	comment=[l]{//},
	commentstyle=\color{purple}\ttfamily,
	frame=lines,
	identifierstyle=\color{black},
	keywordstyle=\color{blue}\bfseries,
	morecomment=[s]{/*}{*/},
	morestring=[b]',
	morestring=[b]",
	ndkeywordstyle=\color{darkgray}\bfseries,
	numbers=left,
	numbersep=8pt,
	numberstyle=\scriptsize,
	sensitive=false,
	showstringspaces=false,
	stepnumber=1,
	stringstyle=\color{red}\ttfamily,
	keywords={
		break,
		case,
		catch,
		catch,
		do,
		else,
		false,
		function,
		if,
		in,
		new,
		null,
		return,
		switch,
		true,
		typeof,
		var,
		while},
	ndkeywords={
		boolean,
		class,
		export,
		implements,
		import,
		this,
		throw}
}

% Listado definido para JSON
% http://tex.stackexchange.com/questions/83085/how-to-improve-listings-display-of-json-files/83100#83100
\lstdefinelanguage{json}{
	backgroundcolor=\color{background},
	basicstyle=\footnotesize,
	breaklines=true,
	captionpos=b,
	frame=lines,
	numbers=left,
	numbersep=8pt,
	numberstyle=\scriptsize,
	showstringspaces=false,
	stepnumber=1,
	literate=
		*{:}{{{\color{punct}{:}}}}{1}
		{,}{{{\color{punct}{,}}}}{1}
	    {\{}{{{\color{delim}{\{}}}}{1}
	    {\}}{{{\color{delim}{\}}}}}{1}
	    {[}{{{\color{delim}{[}}}}{1}
	    {]}{{{\color{delim}{]}}}}{1}
	    {ñ}{{\~{n}}}{1}
}

% Para que las páginas en blanco no tengan cabecera
\makeatletter
\def\clearpage{%
  \ifvmode
    \ifnum \@dbltopnum =\m@ne
      \ifdim \pagetotal <\topskip
        \hbox{}
      \fi
    \fi
  \fi
  \newpage
  \thispagestyle{empty}
  \write\m@ne{}
  \vbox{}
  \penalty -\@Mi
}
\makeatother

\begin{document}
\begin{titlepage}

\newlength{\centeroffset}
\setlength{\centeroffset}{-0.5\oddsidemargin}
\addtolength{\centeroffset}{0.5\evensidemargin}

\noindent\hspace*{\centeroffset}\begin{minipage}{\textwidth}

\centering
\includegraphics[width=0.9\textwidth]{../images/logo_ugr.png}\\[1.4cm]

\textsc{\Large\asunto\\[0.2cm]}
\textsc{\grado}\\[1cm]

{\Huge\bfseries \titulo\\}
\noindent\rule[-1ex]{\textwidth}{3pt}\\[3.5ex]

\centering

\textbf{Autor}\\ {\autor}\\[2.5ex]
\textbf{Profesor}\\ {\profesor}\\[2cm]
\includegraphics[width=0.3\textwidth]{../images/logo_etsiit.png}\\[0.1cm]
\textsc{\escuela}\\
\textsc{---}\\
\ciudad, \today\\
\includegraphics[width=0.3\textwidth]{../images/CC-SA-logo.png}
\end{minipage}
\end{titlepage}
\frontmatter
\begin{center}
{\LARGE\bfseries\titulo}\\
\end{center}
\begin{center}
\autor\
\end{center}

\section*{Resumen}

\bigskip
\noindent{\textbf{Palabras clave}: \textit{\keywords}\\

\subsection*{Los objetivos de esta práctica son:}

\begin{enumerate}

\item
Gestionar	y	desplegar	Máquinas	Virtuales	(MVs) de	forma	automatizada.
\item
Desplegar	 servicios	 y	 aplicaciones	 en	 infraestructuras distribuidas,	 utilizando	
herramientas	de	automatización	de	gestión	y	configuración	de	software.
\item
Diseñar	un	servicio	de	aprovisionamiento	dinámico	de	recursos.
\item
Aprender	a	orquestar	el	conjunto	de	servicios	y	aplicaciones.
\end{enumerate}


\newpage
\thispagestyle{empty}
\
\vspace{3cm}

\noindent\rule[-1ex]{\textwidth}{2pt}\\[4.5ex]

Yo, \textbf{\autor}, alumno de la titulación \textbf{\grado} de la \textbf{\escuela\ de la \universidad}, autorizo la ubicación de la siguiente copia de mi Trabajo (\textit{\titulo}) en la biblioteca del centro para que pueda ser consultada por las personas que lo deseen.

\bigskip
Además, este mismo trabajo está publicado bajo la licencia \textbf{Creative Commons Attribution-ShareAlike 4.0}, dando permiso para copiarlo y redistribuirlo en cualquier medio o formato, también de adaptarlo de la forma que se quiera, pero todo esto siempre y cuando se reconozca la autoría y se distribuya con la misma licencia que el trabajo original. El documento en formato {\tt LaTeX} se puede encontrar en el siguiente repositorio de {\tt GitHub}: \url{https://github.com/erseco/ugr_cloud_computing/tree/master/openstack/}.

\vspace{4cm}

\noindent Fdo: \autor

\vspace{2cm}

\begin{flushright}
\ciudad, a \today
\end{flushright}



\begingroup
\let\cleardoublepage\clearpage
  \tableofcontents
\endgroup
\thispagestyle{empty}
\
\mainmatter
\chapter{Introducción}

\subsubsection{Para	 el	 diseño	 del	 sistema	 de	 almacenamiento	 de	 alta	 disponibilidad	 en	 Cloud,	 hay	 que	
considerar	los	componentes	siguientes	que	deben	ser	desplegados en	la	arquitectura	que	
se	ha	provisto	dentro	de	OpenStack:}

\begin{enumerate}
\item
Servicio	 de	 alta	 disponibilidad / balanceador	 de	 carga.	 Este	 servicio	 será	 el	
encargado	 de	 enrutar y	 balancear	 las	 peticiones	 de	 los	 clientes	 a	 los	 servicios	 de	
almacenamiento	redundantes	que	se	desplegarán.	 Se	usará	HAproxy o	NGINX.	Se	
recomienda	utilizar	NGINX,	debido	a	su	facilidad	de	uso.
\item
Servicio	 de	 almacenamiento.	 Es	 el	 servicio	 de	 almacenamiento	 y	 gestión	 de	 los	
datos	en	 cloud.	Este	 servicio	 se	encontrará	 replicado	en	al	menos	dos	nodos	de	la	
infraestructura	y	dará	soporte	de	almacenamiento	en	cloud.	Se	usará	la	plataforma	
de	almacenamiento	de	ficheros	owncloud.

\item
Servicio	 de	 auto-escalado	 dinámico.	 Este	 servicio	 permite	 que	 se	 creen	 o	
destruyan	 nuevos	 contenedores	 para	 el	 servicio	 de	 almacenamiento.	 El	 servicio	
debe	ser	programado	en	el	lenguaje	que	prefieras.

\item
Servicios	de	bases	de	datos.	El	servicio	de	almacenamiento	requiere	de	un	sistema	
de	 base	 de	 datos	 para	 poder	 almacenar	 toda	la	información	 de	 usuarios,	 ficheros,	
metadatos,	 etc.	 Para	 ello	 como	 mínimo	 se	 usará	 un	 nodo,	 con	
MariaBD/MySQL/PostgreSQL.

\item
Servicio	 de	 autenticación.	 Este	 servicio	 se	 encargará	 de	 habilitar	 el	 acceso	 a	los	
usuarios	 de	 owncloud en	 el	 sistema	 de	 almacenamiento	 en	 cloud.	 El	 servicio	 de	
directorio	 de	 usuarios	 será	 LDAP y	 podrá	 estar	 compartido	 dentro	 del	 nodo	 de	
Base	de	Datos.

\end{enumerate}
\chapter{Implementación}

Para la realización de esta práctica me propuse orquestar utilizando el paradigma PaC (Platform as Code), intenté definir la creación de máquinas con Vagrant y provisionar los diferentes servicios completamente con ansible.

\bigskip
Tras mucha prueba y error con la API de OpenStack y con la particularidad de que no tener disponible lo que en OpenStack denominan \textit{Magic ips} (ips publicas para cada maquina desplegada) conseguí tanto la creación de máquinas como el provisionamiento con ansible en las mismas. Todo esto cuando la máquina lo permitía ya que hubieron diversos problemas que obligaban al reinicio de la máquina OpenStack.

\bigskip
Para poder conectarme a las distintas máquinas tuve que hacer uso de un proxypass ssh que se conectaba al servidor atcstack y de ahí saltaba a las distintas máquinas. Para facilitar la tarea agregué la configuración necesaria al archivo \texttt{~/.ssh/config}

\begin{lstlisting}[language=bash,caption={fichero ~/.ssh/config}]
Host openstack
  Hostname atcstack.ugr.es
  User CC_XXXXXXX

Host machine_158
  User          ubuntu
  HostName      192.168.0.158
  ProxyCommand ssh CC_74003802@atcstack.ugr.es -W %h:%p

Host machine_159
  User          ubuntu
  HostName      192.168.0.159
  ProxyCommand ssh CC_74003802@atcstack.ugr.es -W %h:%p

Host machine_160
  User          ubuntu
  HostName      192.168.0.160
  ProxyCommand ssh CC_74003802@atcstack.ugr.es -W %h:%p

Host machine_161
  User          ubuntu
  HostName      192.168.0.161
  ProxyCommand ssh CC_74003802@atcstack.ugr.es -W %h:%p
\end{lstlisting}

Aunque el código Vagrant funcionaba bien, las continuas caídas del servicio \texttt{OpenStack} hizo que el resto de la configuración la hiciera en instancias en mi propia cuenta de AWS.

\bigskip
Para poder iniciar sesión entre las distintas máquinas se ha generado un par de claves publica-privada  que podemos encontrar en el directorio \texttt{app\_keys}.

\bigskip
La configuración del servidor de base datos fue bastante trivial, agregando usuarios y asignando los roles necesarios. Se ha optado por \texttt{PostgreSQL} por diferenciar el entorno del típico \texttt{LAMP} (Linux + Apache + MySQL + PHP).

\bigskip
La configuración del servidor \texttt{LDAP} si ha sido mas complejo pero al final se ha resuelto instalando el \textit{wrapper} \texttt{ldapscripts} y configurando \texttt{debconf} para asignarle la contraseña de superusuario que queramos al paquete \texttt{slapd}. 

\bigskip
\texttt{OwnCloud} necesita un directorio para persistencia, por lo que si queremos balancear de forma correcta entre máquina la solución mas sencilla es compartir un directorio por NFS, así que instalamos un servidor NFS y configuramos lo necesario para que los demás servidores puedan montar dicha unidad. El código de OwnCloud realiza una serie de bloqueos en el fichero de configuración y dichos bloqueos (\texttt{flock}) no son soportados por NFS por lo que al final se ha optado por inyectar un fichero \texttt{Config.php} dentro del contenedor que evite dichos bloqueos.


\bigskip
Una vez configurado el servidor de BBDDD pasamos a configurar los dos servidores que correrán los contenedores Docker de OwnCloud, en los mismos se han instalado tanto Docker como docker-compose que nos facilitará levantar y escalar contenedores de forma dinámica.

\bigskip
Se ha utilizado el contenedor `owncloud/server' que podemos encontrar en Docker Hub y que está publicado por los propios desarrolladores de OwnCloud.

\bigskip
Aunque dicho contenedor incorpora el plugin user\_ldap, dicho plugin no está activado de forma automática por lo que hemos tenido que definir la variable de entorno \texttt{OWNCLOUD\_APPS\_ENABLE=user\_ldap} para que active el plugin.

\bigskip
Aparte de activar el plugin debemos configurarlo cosa que por defecto no han contemplado los desarrolladores. Al menos permiten ejecutar scripts al finalizar un despliegue. Para ello basta con colocar los scripts que queramos en el directorio /etc/post\_server.d/ del contenedor. Mirando en la documentación de OwnCloud (y tras mucha mas prueba y error) he desarrollado un sencillo script que haciendo uso del cliente \texttt{occ} establece los valores de configuración necesarios para permitir la conexión al servidor LDAP

\begin{lstlisting}[language=bash,caption={bash version}]
#!/usr/bin/bash
LDAP_SERVER={{database}}
# Enable ldap plugin (better do it in env vars)
# su www-data -c "php occ app:enable user_ldap -n -q --no-ansi
# Enable and configure LDAP and external storage plugin
occ ldap:create-empty-config -n -q --no-ansi
occ ldap:set-config 's01' ldapBase dc=example,dc=com -n -q --no-ansi
occ ldap:set-config 's01' ldapBaseGroups dc=example,dc=com -n -q --no-ansi
occ ldap:set-config 's01' ldapBaseUsers dc=example,dc=com -n -q --no-ansi
occ ldap:set-config 's01' ldapCacheTTL 600 -n -q --no-ansi
occ ldap:set-config 's01' ldapConfigurationActive 1 -n -q --no-ansi
occ ldap:set-config 's01' ldapEmailAttribute mail -n -q --no-ansi
occ ldap:set-config 's01' ldapExperiencedAdmin 0 -n -q --no-ansi
occ ldap:set-config 's01' ldapExpertUsernameAttr uid -n -q --no-ansi
occ ldap:set-config 's01' ldapGroupDisplayName cn -n -q --no-ansi
occ ldap:set-config 's01' ldapGroupFilter objectClass=posixGroup -n -q --no-ansi
occ ldap:set-config 's01' ldapGroupFilterMode 0 -n -q --no-ansi
occ ldap:set-config 's01' ldapGroupMemberAssocAttr uniqueMember -n -q --no-ansi
occ ldap:set-config 's01' ldapHost ${LDAP_SERVER} -n -q --no-ansi
occ ldap:set-config 's01' ldapLoginFilter '(&(|(objectclass=posixAccount))(uid=%uid))' -n -q --no-ansi
occ ldap:set-config 's01' ldapLoginFilterEmail 0 -n -q --no-ansi
occ ldap:set-config 's01' ldapLoginFilterMode 0 -n -q --no-ansi
occ ldap:set-config 's01' ldapLoginFilterUsername 1 -n -q --no-ansi
occ ldap:set-config 's01' ldapNestedGroups 0 -n -q --no-ansi
occ ldap:set-config 's01' ldapPagingSize 500 -n -q --no-ansi
occ ldap:set-config 's01' ldapPort 389 -n -q --no-ansi
occ ldap:set-config 's01' ldapQuotaAttribute mailQuota -n -q --no-ansi
occ ldap:set-config 's01' ldapTLS 0 -n -q --no-ansi
occ ldap:set-config 's01' ldapUserDisplayName cn -n -q --no-ansi
occ ldap:set-config 's01' ldapUserFilter objectClass=posixAccount -n -q --no-ansi
occ ldap:set-config 's01' ldapUserFilterMode 0 -n -q --no-ansi
occ ldap:set-config 's01' ldapUuidGroupAttribute auto -n -q --no-ansi
occ ldap:set-config 's01' ldapUuidUserAttribute auto -n -q --no-ansi
\end{lstlisting}

Una vez configurado LDAP de forma automática en nuestro contenedor procedimos a levantar mas de un contenedor. OwnCloud por defecto no contempla esta configuración por lo que tuvimos que hacer un bypass del script \texttt{/usr/local/bin/owncloud-installed}, que es uno de los scripts de comprobación que ejecuta a la hora de inicializar el contenedor y que crea la base de datos. En nuestro caso lo cambiamos para que en caso de existir la base de datos retornara un valor para que pensara que estamos ante una actualización en lugar de en una instalación.

\begin{lstlisting}[language=bash,caption={comprobacion de que no está previamente configurado OwnCloud}]
#!/usr/bin/env bash
set -e
if [[ $(PGPASSWORD=$OWNCLOUD_DB_PASSWORD psql -U $OWNCLOUD_DB_USERNAME -h {{database}} $OWNCLOUD_DB_NAME -tAc "SELECT to_regclass('oc_users');") != "oc_users" ]]
then
  exit 1
else
  exit 0
fi
\end{lstlisting}


Una vez resuelto este problema ya hemos podido levantar múltiples contenedores con la sencilla instrucción \texttt{docker-compose up --scale owncloud=N} siendo N el numero de contenedores que queremos tener. El propio docker-compose se encarga de reescalar si el numero es distinto del numero de contenedores en ejecución.

\bigskip
Hay que tener en cuenta que al permitir levantar contenedores de forma dinámica no podemos fijar los puertos, así que en el docker-compose le indicamos que asigne un puerto aleatorio a cada contenedor que redirija al puerto 80 que es el que está sirviendo OwnCloud.

\bigskip
Una vez configuradas las máquinas Docker hemos pasado a configurar el balanceador, que consiste en un nginx  configurado para incluir un fichero de configuración donde sacar la lista de servidores que llamará a través de \texttt{proxypass}.

\bigskip
Dicha lista de contenedores la hemos extraído de la salida del comando \texttt{docker ps --format "{{.Ports}}"} que nos devuelve los puertos asignados a los contenedores. De esta forma nuestro sistema permite levantar tantos contenedores como queramos sin tener que fijar un numero de puertos.

\bigskip
Una vez configurado el script de balanceo, hemos procedido a probar que tal se comportaba viendo que como sospechábamos \texttt{OwnCloud} no está preparado de forma nativa para ser balanceado al guardar las sesiones de los usuarios en el sistema interno de \texttt{PHP} que guarda en disco las mismas siendo lo recomendado en estos casos es guardar la configuración de las sesiones en la base de datos.

\bigskip
Para el balanceo de carga se ha utilizado un sencillo script en bash que incorpora las funciones necesarias para consultar el estado de \texttt{NGINX} y conectarse a los dos nodos \texttt{Docker} para mediante docker-compose hacer un reescalado automático. El script detecta automáticamente los puertos en los que se ha levantado Docker y los incorpora a la lista del \textit{upstream} de \texttt{NGINX} haciendo un reinicio controlado (\textit{reload}) que permite que el servicio no esté caído en ningún momento como si pasaría con un reinicio normal (\textit{restart}). Dicho scripts guarda un log de las tareas realizadas en el archivo \texttt{saed.log}


\bigskip
Para dicho balanceo nos hemos limitado a realizar la operación matemática proporcionada por el profesor que en nuestras pruebas de carga con \texttt{siege} no nos ha dado valores diferentes de 1 o 2, quizá por filtrado de la propia AWS, lo que hemos hecho es utilizar ese numero como valor del numero de contenedores a levantar en cada nodo.






\chapter{Conclusiones}

Esta práctica ha supuesto un gran reto ya que en mi opinión el sistema \texttt{OwnCloud} no está preparado para ser balanceado de forma sencilla, al menos utilizando los contenedores \texttt{Docker} proporcionados por los desarrolladores. De igual forma dichos contenedores no están preparados para integrar LDAP sin realizar diversas tareas de programación, de igual forma el código de OwnCloud no está utilizando correctamente las sesiones PHP para funcionar en un entorno balanceado.

\bigskip
Aun así me ha permitido trabajar con \texttt{OpenStack} así como jugar con una configuración que requería de ingenio para poder saltarse las restricciones de la plataforma.
\chapter{Bibliografía}

Se ha utilizado la documentación oficial de los proyectos, así como el código fuente de los proyectos utilizados.

\begin{itemize}
\item{\tt Documentación de Docker}: \url{https://docs.docker.com/compose/environment-variables/}
\item{\tt Documentación de Nginx}: \url{https://nginx.org/en/docs/}
\item{\tt Documentación de OwnCloud}: \url{https://doc.owncloud.org/}
\item{\tt Documentación de Ansible}: \url{http://docs.ansible.com/}
\item{\tt StackOverflow}: \url{https://stackoverflow.com}
\item{\tt Creative Commons Share Alike 4.0}: \url{https://creativecommons.org/licenses/by-sa/4.0/}
\end{itemize}




\newpage \
\thispagestyle{empty}
\end{document}